\subsection{A ruled residual and its subtraction cost}
\label{subsec:ruled-residual-gauduchon-cost}

On the ruled surface of \cref{sec:nonclosed-ruled-trace-masses}, the Gauduchon variational volume has a different dependence on the data from the two volumes of an ordinary type. Keep $Y=\mathbb P^1\times\mathbb P^1$, $\theta=cw_1+aw_2$, $0\leq b<c$ and $\beta=(c-b)w_1+aw_2$. Let $m>0$ be smooth on the first curve, with $\int mw=1$, and put
\[
 \mu=mw_1\wedge w_2,\qquad k=c-b,\qquad H=\int_{\mathbb P^1}\frac{m}{a}w.
\]
In dimension two the Gauduchon tests are all smooth positive $(1,1)$-forms $T$ with $\ddc T=0$, and $D_\mu(T)=\int_Y\sqrt{T^2/\mu}\,\mu$.

\enlargethispage{4\baselineskip}
\begin{proposition}\label{prop:ruled-residual-gauduchon-volume}\index{Gauduchon variational volume!ruled residual}
The all-test infimum is attained and
\begin{equation}\label{eq:ruled-residual-gauduchon-volume}
 \mathcal V_\mu(\theta;bD)=\mathcal V_\mu(\beta)=\frac{2(c-b)}{H},
 \qquad \mathcal V_\mu(\theta)=\frac{2c}{H}.
\end{equation}
Writing $V=2k/H$ and $g=Vm/(2a)$, a minimizing test is the closed positive form
\[
 T=\frac ga w_1+w_2.
\]
\end{proposition}
\begin{proof}
Since $\int gw=k$, solve $\ddc f=(g-k)w$ on the first sphere. Then $\Omega=\beta+\ddc f=gw_1+aw_2>0$ and $\Omega^2=V\mu$. For every Gauduchon test, Stokes and the matrix arithmetic--geometric mean inequality give
\[
 \int_Y\beta\wedge T=\int_Y\Omega\wedge T
       \geq\sqrt V\,D_\mu(T).
\]
The displayed $T$ is closed, since $g/a$ depends only on the first factor, and $T=a^{-1}\Omega$, so equality holds pointwise. Explicitly,
\[
 \int_Y\beta\wedge T=2k,
 \qquad D_\mu(T)=\sqrt V\int\frac ma w=\sqrt V H.
\]
This proves the infimum over all tests. The curvature of $D$ is $w_1$, so the divisorial residual is indeed $\beta$. Replacing $k$ by $c$ gives the formula for $\theta$.
\end{proof}

Fix a real constant $d$ and a smooth real first-factor function $e$. On an interval where $c_s=c-sd>b$ and $a_s=a-se>0$, keep $b,m$ fixed and put
\[
 k_s=c_s-b,\qquad H_s=\int\frac{m}{a_s}w,
 \qquad V_s=\frac{2k_s}{H_s},\qquad
 T_s=\frac{V_sm}{2a_s^2}w_1+w_2.
\]
\begin{corollary}\label{cor:ruled-residual-subtraction-cost}
For the residual background $\beta_s=k_sw_1+a_sw_2$,
\begin{equation}\label{eq:ruled-residual-subtraction-cost}
 \frac{\mathrm d}{\mathrm ds}\sqrt{V_s}
     =-\frac{\int_Y(d w_1+e w_2)\wedge T_s}{D_\mu(T_s)}.
\end{equation}
Consequently, for every compact segment $[s_0,s_1]$ in that interval,
\[
 \sqrt{V_{s_0}}-\sqrt{V_{s_1}}
   =\int_{s_0}^{s_1}
       \frac{\int_Y(d w_1+e w_2)\wedge T_s}{D_\mu(T_s)}\,\mathrm ds.
\]
\end{corollary}
\begin{proof}
We have $H_s'=\int em/a_s^2\,w$, $D_\mu(T_s)=\sqrt{V_s}H_s$ and
\[
 \int_Y(d w_1+e w_2)\wedge T_s=d+\frac{V_s}{2}H_s'.
\]
Differentiating $V_s=2k_s/H_s$ gives \eqref{eq:ruled-residual-subtraction-cost}; integration proves the second identity. This is the explicit first variation of \cref{prop:gaud-volume-first-variation} along the residual family.
\end{proof}

For $c=1$, $b=1/4$, $a=2+\tau$ and $m=1$, the height $\tau$ has distribution $\mathrm d\tau/2$ under $w$. Hence $H=(\log3)/2$, and the residual and original variational volumes are $3/\log3$ and $4/\log3$. The ordinary-type volumes from \eqref{eq:nonclosed-ruled-ordinary-volumes} are instead $3/2$ and $9/2$, independent of $m$. The subtraction identity differentiates the density-dependent variational volume. Pure-pullback invariance of the ordinary-type masses does not identify an arbitrary smooth positive probability density upstairs with the degenerate pullback of $\mu$.
